\section{방법}
\label{sec:method}

\begin{figure*}[t]
\centering
\includegraphics[width=\textwidth]{figures/fig_scopead_overview.pdf}
\caption{\textbf{Overview of \method{}.} 
(a) 정상 서포트 표본으로부터 leave-source-out (1-shot의 경우 leave-token-out) 교차 매칭을 통해 48개의 독립 메모리 뱅크를 구축하고 참조 척도 $c_m$을 추정한다. 
(b) 코사인 1-NN 거리 $d_m$을 해당 척도로 정규화하여 공통 단위 점수 $z_m = d_m / c_m$을 도출한다. 
(c) 전역 문맥을 담은 깊은 계층을 주축으로 얕은 계층의 단서를 보조 결합하는 비대칭 선형 가중합을 수행한다. 
(d) 배그 평균 및 원본 해상도 보간 후, 각 픽셀 위치에서 정렬된 해상도 점수에 Rényi 순서 통계량 가중치를 적용하여 최종 이상 맵을 집계한다. 가중치는 해상도 크기가 아닌 점수 크기 순위에 따라 비대칭적으로 부여된다. 예시는 결함 픽셀 $p$의 실제 점수다(평균 1.41).}
\label{fig:method}
\end{figure*}

\method는 관측 확장(observation expansion)과 정상 참조 기반 점수 통합을 결합하여, 소수 서포트와 환경 변화로 인한 High-Score Ambiguity를 완화하는 이상 국소화 프레임워크다(\cref{fig:method}).
그 중심에서 \method는 표현 계층과 입력 해상도를 독립 매칭 단위로 분리하는 관측 확장을 수행하여, 특정 계층이나 해상도가 거리 계산을 지배하는 간섭을 배제하고 단일 영역에 대한 다각적 시각 단서를 확보한다.
이렇게 확보된 다관측 거리들을 효과적으로 결합하기 위해 세 가지 요소를 단계적으로 통합한다:
(1) 서포트 교차 매칭 분포의 상위 분위수로 기준 단위를 정렬하는 \textbf{정상 참조 척도화(Reference Scaling)},
(2) 전역 문맥을 담은 심층 특징을 기저로 삼고 얕은 계층의 반응을 수축 결합하는 \textbf{계층별 비대칭 결합(Asymmetric Layer Fusion)},
그리고 (3) 위치별 해상도 점수 순위에 따라 상위 반응을 안정적으로 보존하는 \textbf{Rényi 순서 통계량 집계(Rényi Order-Rank Aggregation)}이다.
추론 시 \method는 백본의 추가 학습이나 쿼리 적응 없이, 정상 서포트의 무감독 통계량만을 활용하여 단일 패스에서 픽셀 단위 결함 맵과 이미지 점수를 산출한다.

\subsection{문제 설정}
\label{sec:method:overview}

입력으로 $N$장의 정상 서포트 이미지 집합 $\mathcal{S} = \{x_1^s, \dots, x_N^s\}$, 테스트 쿼리 이미지 $x^q$, 그리고 동결된 사전학습 비전 트랜스포머(ViT) 백본이 주어진다.
특징을 추출할 백본 계층 집합을 $\mathcal{L}$, 입력 해상도(짧은 변 기준) 집합을 $\mathcal{R}$, 서브샘플 배그(bag) 수를 $B$로 정의한다.
본 연구에서는 $\mathcal{L}=\{5, 11, 23\}$, $\mathcal{R}=\{448, 672, 896, 1120\}$, $B=4$를 사용하여 총 $|\mathcal{L}| \times |\mathcal{R}| \times B = 48$개의 독립 관측을 구성한다(\cref{sec:experiments:setup}).
\method는 백본의 추가 학습이나 쿼리 적응(query adaptation) 없이, 오직 정상 서포트의 무감독 통계량만을 기반으로 픽셀 단위 결함 국소화 맵 $A(x^q) \in \mathbb{R}^{H \times W}$ 및 이미지 단위 이상 점수 $s(x^q) \in \mathbb{R}$를 산출한다.

\subsection{관측 확장과 독립 참조 매칭}
\label{sec:method:expansion}

\paragraph{관측별 독립 매칭 (Independent Matching).}
획득 조건이 달라지면 정상 영역이라도 특정 계층의 국소 질감이나 특정 해상도의 패치에서 거리가 급격히 상승할 수 있다.
기존처럼 다층 특징을 단일 벡터로 연결해 매칭하면~\cite{patchcore}, 분산이나 척도가 큰 특정 관측이 최근접 이웃 탐색 전체를 지배하여 정상 외관 변화에 따른 국소적 거리 상승이 다른 모든 계층의 비교 단서를 왜곡하게 된다.
이러한 간섭을 차단하기 위해 \method는 계층 $l \in \mathcal{L}$, 해상도 $r \in \mathcal{R}$, 배그 인덱스 $b \in \{1, \dots, B\}$의 조합 $m = (l, r, b)$를 개별 관측 단위로 분리하고, 각 관측마다 독립적인 정상 메모리 뱅크 $B_m$을 구축한다.
쿼리 이미지 $x^q$의 공간 위치 $p$에서 추출된 $\ell_2$ 정규화 패치 토큰 $f_m(x^q, p)$는 오직 대응하는 뱅크 $B_m$ 내에서만 독립적으로 최근접 이웃을 탐색하여 코사인 거리 점수를 계산한다:
\begin{equation}
  d_m(x^q, p) = \min_{u \in B_m} \left[ 1 - f_m(x^q, p)^\top u \right].
  \label{eq:cosine-dist}
\end{equation}
이 분리 매칭 구조에서는 각 관측이 고유한 표현·공간 척도 안에서 가장 가까운 정상을 개별적으로 찾는다.
그 결과 특정 척도에서 발생한 정상 외관 변동의 섭동이 다른 관측으로 전파되지 않고 격리되며, 국소 질감 단서와 전역 문맥 단서가 독립적인 거리 신호로 온전히 보존되어 후속 통합 단계에서 상호 검증의 기반을 마련한다.

\paragraph{패치 서브샘플 배깅 (Patch Subsample Bagging).}
입력 해상도가 증가함에 따라 서포트 패치 토큰의 수가 급격히 팽창하며, 단일 무작위 샘플링은 토큰 구성에 따른 추정 분산을 유발할 수 있다.
이를 완화하기 위해 각 해상도 $r$에서 서포트 이미지당 $K$개($K=256$)의 패치 토큰을 비복원 균등 추출하여 단일 배그를 구성하며, 서로 다른 난수 시드로 $B$개의 배그를 생성한다.
이때 일관된 공간 대응을 유지하기 위해 동일 배그 내의 모든 계층은 완전히 일치하는 공간 좌표의 토큰을 공유한다. 각 배그에서 산출된 이상 점수는 이후 \cref{eq:bag-mean}에서 앙상블 평균된다.

\subsection{정상 참조 기반 비대칭 선형 결합}
\label{sec:method:fusion}

독립적으로 획득한 원시 거리 $d_m$을 통합하기 위해서는 (1) 관측 간 척도 불균형(scale discrepancy)과 (2) 계층별 의미론적 민감도 차이를 단계적으로 다루어야 한다.
특징 공간의 차원과 표현 특성에 따라 원시 거리의 통계적 척도가 상이하므로(\cref{fig:scale-calib}), 보정 없이 결합하면 스케일이 큰 관측의 오경보가 결합 점수에 그대로 반영된다.
또한 척도를 정규화하더라도 국소 질감에 민감한 얕은 계층은 조명이나 자세 변화 등 정상 외관 변동에 큰 거리를 유발하기 쉽다.
따라서 \method는 서포트 교차 매칭 기반의 척도 $c_m$으로 단위를 정렬한 뒤, 비대칭 계층 가중치 $\alpha_l$을 적용하여 선형 결합을 수행한다.

\begin{figure}[t]
\centering
\includegraphics[width=\linewidth]{figures/fig_scopead_calibration_paper.pdf}
\caption{\textbf{정상 참조 보정(Q95)을 통한 관측 간 거리 척도 정렬.}
(a) Vial 2-shot 정상 쿼리에서 12개 관측(4개 해상도 $\times$ 3개 계층)의 상대 정상 척도. 보정 전(회색)에는 관측마다 정상 거리 척도가 최대 1.60배까지 불균일하여 특정 관측이 결합을 왜곡하지만, Q95 보정 후(청색)에는 모든 관측이 공통 기준선 근처로 균일하게 정렬된다(최대 편차 1.60배 $\rightarrow$ 1.24배).
(b) MVTec AD 2 전 클래스(8개 클래스 $\times$ 2개 shot)에서 관측 간 척도 불균형(최대/최소 척도 비율). 16개 조건 중 14개 조건에서 보정 후 척도 편차가 대각선($y=x$) 아래로 크게 감소하여 안정적으로 정렬됨을 보인다.}
\label{fig:scale-calib}
\end{figure}

\paragraph{서포트 교차 매칭 기반 척도 추정 (Cross-Validation Scale Estimation).}
미지의 쿼리 이미지와의 매칭 환경을 모사하기 위해, 서포트 집합 내부의 교차 매칭 거리 표본 $D_m$을 활용하여 관측별 기준 척도를 추정한다.
다중 샷($N > 1$) 환경에서는 동일 이미지 내의 자기 유사 패치와 매칭되어 거리가 과소 추정되는 편향을 방지하고자, 서포트 이미지 $i$의 토큰 부분집합 $B_m^{(i)}$를 해당 이미지를 제외한 나머지 뱅크 토큰들과 매칭하는 Leave-Source-Out (LSO) 방식을 적용한다:
\begin{equation}
  D_m^{\text{LSO}} = \bigcup_{i=1}^N \left\{ \min_{u \in B_m \setminus B_m^{(i)}} \left[ 1 - v^\top u \right] : v \in B_m^{(i)} \right\}.
  \label{eq:lso}
\end{equation}
단일 샷($N = 1$) 환경에서는 이미지 간 교차 비교가 불가능하므로, 단일 이미지 내에서 자기 자신을 제외한 타 위치 토큰들과 매칭하는 Leave-Token-Out (LTO) 표본을 대리 통계량으로 사용한다:
\begin{equation}
  D_m^{\text{LTO}} = \left\{ \min_{j \neq i} \left( 1 - v_i^\top v_j \right) : i = 1, \dots, K \right\}.
  \label{eq:lto}
\end{equation}

각 관측의 정상 거리 분포는 비대칭성과 꼬리 두께가 상이하므로, 특정 확률분포 형태를 가정하는 모수적 통계량 대신 비모수적 순서 통계량인 상위 분위수 $\mathcal{Q}_q$로 기준 척도 $c_m$을 정의하고 정규화 점수 $z_m$을 산출한다:
\begin{equation}
  c_m = \max \left\{ \mathcal{Q}_q(D_m), \, \epsilon \right\}, \qquad z_m(x^q, p) = \frac{d_m(x^q, p)}{c_m},
  \label{eq:scale-cm}
\end{equation}
여기서 $q=0.95$는 서포트 집합 내에서 관측되는 정상 거리의 상위 경계를 대표하는 분위수이며, $\epsilon=10^{-4}$는 수치적 안정성을 위한 하한값이다.
기준 척도 $c_m$은 각 관측 뱅크가 포착한 정상 외관 변동의 실질적 상한선을 뜻한다.
따라서 $z_m = d_m / c_m$ 변환은 이질적인 원시 거리들을 '정상 변동 상한 대비 초과율'이라는 동일한 무차원 판정선($z_m = 1.0$) 위에 정렬한다(\cref{fig:scale-calib}).
이를 통해 원시 거리의 절대적인 기저 크기가 큰 특정 계층이나 해상도가 결합 점수를 일방적으로 지배하여, 정상 외관 변동이 결합 점수에서 오경보로 전파되는 왜곡을 완화한다.

\paragraph{계층 간 비대칭 결합 (Asymmetric Layer Fusion).}
척도가 정렬된 점수 $z_m$ 공간에서도 계층 간 역할 차이가 존재한다.
조명이나 표면 텍스처 등 획득 조건이 달라지면 국소 외관에 반응하는 얕은 계층은 정상 영역에서도 거리가 크게 치솟기 쉽다.
반면 전역 문맥을 담은 심층 계층은 획득 환경이 변화하더라도 정상 구조의 의미론적 일관성을 유지하므로, 정상 영역의 점수를 상대적으로 낮게 보존하는 안정적인 기저 판정선을 제공한다.
이를 반영하여 깊은 계층들($\mathcal{L}_{\text{deep}}=\{11, 23\}$)의 정규화 점수 평균을 기저 점수로 삼고, 정상 변동에 취약한 얕은 계층($l=5$)은 작은 가중치 $\lambda=0.125$로 축소 결합한다:
\begin{equation}
  S_{r,b}(x^q, p) = \frac{1}{2} \sum_{l \in \{11, 23\}} z_{l,r,b}(x^q, p) + \lambda \, z_{5,r,b}(x^q, p).
  \label{eq:layer-fusion}
\end{equation}
이 비대칭 결합은 환경 변화로 인한 얕은 계층의 오경보 유입을 전역 문맥 기저선으로 억제하면서, 국소 미세 결함 단서만을 선별적으로 보완한다.
정규화 점수 정의($z_m = d_m/c_m$)에 따라, 이는 원시 거리 $d_{l,r,b}$를 계층 가중치 $\alpha = (0.125, 0.5, 0.5)$와 각 관측의 기준 척도 $c_{l,r,b}$의 비율인 $w_{l,r,b} = \alpha_l / c_{l,r,b}$로 선형 결합하는 것과 같다.

\subsection{다해상도 순서 통계량 집계}
\label{sec:method:aggregation}

앞선 계층 결합을 거친 이상 점수는 여전히 $R$개의 입력 해상도별 맵으로 분리되어 있다.
먼저 각 해상도 $r$에서 $B$개 배그의 점수를 평균하여 패치 서브샘플링에 따른 추정 분산을 완화한 뒤, 양선형 보간(bilinear interpolation)을 통해 모든 해상도 맵 $S_r$을 원본 이미지 크기 $(H, W)$로 공간 정렬한다:
\begin{equation}
  S_r(x^q, p) = \frac{1}{B} \sum_{b=1}^B S_{r,b}(x^q, p).
  \label{eq:bag-mean}
\end{equation}

\paragraph{Rényi 순서 통계량 집계 (Order-Statistics Aggregation).} 
공간 축이 정렬된 다해상도 점수 맵들을 결합할 때, 단순 산술 평균($1/R=0.25$)은 특정 해상도에서만 뚜렷하게 관측되는 유효 결함 피크를 타 해상도의 낮은 점수와 섞어 희석시킨다.
반대로 결함 피크를 보존하기 위해 단순히 최댓값($\max$)만을 취하면, 정상 영역에서 단일 해상도에 우연히 발생한 단발성 잡음 스파이크가 그대로 최종 결함 점수로 채택되어 오경보를 유발한다.
따라서 다해상도 결합에는 최상위 반응의 피크는 온전히 보존하면서도, 단일 관측 잡음의 돌출을 다관측 지지 반응으로 완충(damping)하는 비대칭 가중 체계가 필요하다.

이러한 상위 편향 가중치를 자의적인 휴리스틱 튜닝 없이 결정하기 위해, 이상 탐지의 극치 꼬리 거동(tail behavior)을 다루는 표준 지수 분포의 순서 통계량 기댓값 간격(Rényi representation, $a_k = \sum_{j=1}^k \frac{1}{R-j+1}$)을 참조 모형으로 채택한다.
조화 수열(harmonic sum)에 기반한 이 기댓값 구조는 최상위 순위로 갈수록 가중치가 자연스럽게 완만히 증가하는 특성을 갖는다.
$R=4$일 때 정규화된 Rényi 순위 가중치 $w = (\frac{1}{16}, \frac{7}{48}, \frac{13}{48}, \frac{25}{48}) \approx (0.063, 0.146, 0.271, 0.521)$를 통해 다해상도 결합 맵을 산출한다:
\begin{equation}
  \tilde{A}(x^q, p) = \sum_{k=1}^R w_k \, S_{(k)}(x^q, p).
  \label{eq:scale-renyi-mean}
\end{equation}
이 가중 체계는 다해상도 이상 맵 결합에 두 가지 실질적인 이점을 제공한다:
(1) 최상위 점수 순위에 약 52.1\%($25/48$)의 지배적인 가중치를 부여하여 특정 해상도에 집중된 결함 피크 신호의 희석을 직접적으로 방어한다.
(2) 동시에 차상위 순위들($k=2, 3$)에도 약 41.7\%($20/48$)의 완만한 지지 가중치를 분배함으로써, $\max$ 집계처럼 한 해상도에서만 외톨이로 튄 국소 잡음 스파이크가 최종 결함으로 과대 채택되는 불안정성을 효과적으로 완충한다.
특히 앞 절의 \textbf{정상 참조 척도 보정($c_m$)}이 모든 해상도의 정상 기저 점수를 $z=1.0$ 기준으로 균일화해 두었기 때문에, 특정 해상도의 기저 잡음이 상위 순위를 차지해 오경보를 유발하는 위험 없이 유효 결함 반응만을 안전하게 보존할 수 있다.
Rényi 집계와 보정의 상호작용 메커니즘 분석은 \cref{sec:experiments:mechanism}에서 다룬다.

\paragraph{이미지 수준 이상 점수.}
최종 픽셀 단위 결함 국소화 맵은 $A(x^q) = \tilde{A}(x^q)$로 정의된다.
이미지 수준 이상 점수 $s(x^q)$는 $A(x^q)$ 내 상위 1\% 픽셀 집합 $\Omega_{0.01}$($|\Omega_{0.01}| = \max\{1, \lfloor 0.01HW \rfloor\}$)의 평균값으로 정의된다:
\begin{equation}
  s(x^q) = \frac{1}{|\Omega_{0.01}|} \sum_{p \in \Omega_{0.01}} A(x^q, p).
  \label{eq:image-score}
\end{equation}
